\section{Mean Bias as the Primary Source of Second-Moment Inflation}
\label{subsec:mean_second_moment}

We now turn to the optimizer-level consequence of activation bias.
While the previous subsection established that mean bias dominates
elementwise outliers, adaptive optimizers operate on squared gradients.
It is therefore essential to understand how activation bias propagates
into second-moment statistics.

\paragraph{Gradient outer-product structure.}

For a linear layer with weight matrix $W \in \mathbb{R}^{H_{\text{out}}\times H_{\text{in}}}$,
the per-token gradient is the outer product

\begin{equation}
g = \delta x^\top,
\end{equation}

where $x \in \mathbb{R}^{H_{\text{in}}}$ is the forward activation vector
and $\delta \in \mathbb{R}^{H_{\text{out}}}$ is the gradient output vector.

Adam and related adaptive optimizers maintain the elementwise second moment $v = \mathbb{E}[g \odot g]$.
To rigorously track dimensions, consider the weight connecting input coordinate $j$ to output coordinate $i$. The second moment for this parameter is

\begin{equation}
v_{ij} = \mathbb{E}[g_{ij}^2]
   = \mathbb{E}[\delta_i^2 x_j^2].
\end{equation}

\paragraph{Activation decomposition.}

Decompose the activation coordinate $x_j$ into its expected mean and centered fluctuation:

\begin{equation}
x_j = \mu_j + \tilde{x}_j,
\qquad
\mathbb{E}[\tilde{x}_j] = 0.
\end{equation}

Substituting into $v_{ij}$ yields

\begin{equation}
v_{ij}
=
\mathbb{E}[\delta_i^2 (\mu_j + \tilde{x}_j)^2]
=
\mathbb{E}[\delta_i^2 \mu_j^2]
+
\mathbb{E}[\delta_i^2 \tilde{x}_j^2]
+
2\mu_j \mathbb{E}[\delta_i^2 \tilde{x}_j].
\end{equation}

Crucially, because the backward gradient energy $\delta_i^2$ is empirically found to be approximately uncorrelated
with the forward activation fluctuation $\tilde{x}_j$
(Section~\ref{subsec:mean_emergence}),
the expectation of their product separates. Since the fluctuation is zero-mean ($\mathbb{E}[\tilde{x}_j] = 0$), the cross term vanishes in expectation:

\begin{equation}
2\mu_j \mathbb{E}[\delta_i^2 \tilde{x}_j]
\approx
2\mu_j \mathbb{E}[\delta_i^2]\mathbb{E}[\tilde{x}_j]
= 0.
\end{equation}

Thus the second moment separates cleanly into two additive components:

\begin{equation}
v_{ij}
\approx
\underbrace{\mathbb{E}[\delta_i^2]\mu_j^2}_{\text{bias-driven term}}
+
\underbrace{\mathbb{E}[\delta_i^2 \tilde{x}_j^2]}_{\text{variance-driven term}}.
\label{eq:v_decomposition}
\end{equation}

\paragraph{High-dimensional amplification.}

Let $\bar{\mu}$ denote the average coordinate magnitude of the mean vector.
In $H_{\text{in}}$ dimensions,

\begin{equation}
\|\mu\|_2^2 = H_{\text{in}} \bar{\mu}^2.
\end{equation}

Even when $\bar{\mu}$ is moderate,
the coherent alignment across $H_{\text{in}}$ coordinates
creates a dense $\mu_j^2$ term that contributes uniformly
to all output dimensions $i$ connected to that input feature $j$.
In contrast, $\tilde{x}_j$ is approximately isotropic,
so its squared contribution spreads across directions
without coherent alignment.
Consequently, for extreme percentiles of the matrix $v$
(e.g., $p99.9$ or $\max$),
the bias-driven term in \eqref{eq:v_decomposition}
dominates.

\paragraph{Extreme-value scaling.}

From \eqref{eq:v_decomposition}, if the structural mean magnitude heavily outweighs the variance, $|\mu_j| \gg \sqrt{\mathbb{E}[\tilde{x}_j^2]}$,
then

\begin{equation}
v_{ij} \sim \mathbb{E}[\delta_i^2] \mu_j^2,
\end{equation}

and extreme values scale quadratically with the coherent mean.
This mathematically guarantees why activation bias directly inflates
the heavy tail of $v$.

\paragraph{Empirical verification.}

We quantify this effect across shallow and deep layers
and across early (10k) and late (180k) training stages.
Table~\ref{tab:adam_v_bias} reports the $p99.9$ and maximum values
of $v$ before and after mean removal. Across layers and training stages,
mean removal consistently contracts the extreme tail of $v$.
In early layers and late training,
reductions exceed $4\times$.
Whitening of the residual spike provides additional but smaller gains.


\begin{table}[t]
\centering
\caption{Effect of mean removal on extreme second-moment statistics.}
\label{tab:adam_v_bias}
\begin{tabular}{lcccc}
\toprule
Layer & Step & $p99.9(v)$ (raw) & $p99.9(v)$ (mean removed) & Reduction \\
\midrule
3 & 10k  & $1.60\times10^{-10}$ & $3.44\times10^{-11}$ & $4.6\times$ \\
3 & 180k & $1.00\times10^{-10}$ & $2.50\times10^{-11}$ & $4.0\times$ \\
21 & 10k & $1.67\times10^{-11}$ & $1.59\times10^{-11}$ & $1.05\times$ \\
21 & 180k& $9.23\times10^{-11}$ & $2.21\times10^{-11}$ & $4.2\times$ \\
\bottomrule
\end{tabular}
\end{table}


These results demonstrate that the dominant driver of extreme second-moment statistics is the rank-one mean component. Residual spike variance contributes secondary structure. The long-tail variance contributes negligibly to extreme $v_{ij}$.

Since adaptive learning-rate scaling depends on $1/\sqrt{v_{ij}}$,
bias-driven inflation suppresses updates along semantically meaningful
directions.
Removing the mean directly eliminates the first-order source of
second-moment distortion,
thereby restoring balanced effective learning rates.


In summary, activation mean bias is the primary source of second-moment inflation.
Its effect is quadratic, dense, and amplified by dimensionality.
Mean removal therefore targets the first-order cause
of optimizer instability.